% Owner: Vinh
\subsection{Q5: Convolution}

% Vấn đề giải quyết
\subsubsection{Vấn đề giải quyết}
Q5 tiếp tục bài toán phân loại chữ số viết tay MNIST của Q3: mỗi ảnh xám
kích thước $28 \times 28$ cần được xếp vào một trong 10 lớp (chữ số 0 đến
9). Điểm khác biệt là trước khi đưa vào mạng nhiều lớp, ảnh được đi qua một
lớp tích chập (convolution) do sinh viên tự cài đặt trong hàm
\texttt{Convolve}, chỉ dùng các phép toán tensor cơ bản, không dùng hàm
\texttt{conv2d} có sẵn của PyTorch.

Đầu vào là vector $x \in \mathbb{R}^{784}$, được định dạng lại thành ma trận
ảnh $X \in \mathbb{R}^{28 \times 28}$. Đầu ra là vector logits
$z \in \mathbb{R}^{10}$, và lớp dự đoán là $\hat{c} = \arg\max_k z_k$.

Lớp tích chập khắc phục hai hạn chế của mô hình MLP ở Q3, vốn duỗi ảnh
thành vector:
\begin{itemize}
  \item \emph{Bỏ qua vị trí điểm ảnh:} hai điểm ảnh kề nhau chỉ là hai đầu
  vào độc lập. Nếu hoán vị thứ tự các điểm ảnh (giống nhau cho mọi ảnh), MLP
  vẫn học được như cũ, tức là nó không dùng thông tin về cấu trúc không gian.
  \item \emph{Không chia sẻ trọng số:} một nét dọc ở góc trên trái và cùng
  nét đó ở giữa ảnh được xử lý bởi các trọng số khác nhau, nên mô hình phải
  học lại cùng một mẫu tại mỗi vị trí.
\end{itemize}

Autograder chấm hai phần: hàm \texttt{Convolve} phải cho kết quả khớp với
phép tích chập tham chiếu trên các ma trận ngẫu nhiên (1 điểm), và mô hình
phải đạt độ chính xác tối thiểu 80\% trên tập test (2 điểm). Dữ liệu của Q5
nhỏ hơn nhiều so với Q3: 200 ảnh huấn luyện, tập validation và tập test mỗi
tập 50 ảnh. Cả hai tập này đều được lấy từ 100 ảnh đầu tiên của chính tập
huấn luyện; ảnh hưởng của điều này được phân tích ở phần Kết quả.

% Kiến thức sử dụng
\subsubsection{Kiến thức sử dụng}

\textbf{Phép tích chập 2 chiều.} Cho ảnh $X \in \mathbb{R}^{H \times W}$ và
kernel (bộ lọc) $K \in \mathbb{R}^{k_H \times k_W}$. Kernel được đặt lên một
vùng con (patch) cùng kích thước của ảnh; tại mỗi vị trí, nhân từng phần tử
của patch với phần tử tương ứng của kernel rồi cộng lại, được một giá trị
của bản đồ đặc trưng (feature map) $Y$:
\begin{equation}
Y_{i,j} = \sum_{a=0}^{k_H-1} \sum_{b=0}^{k_W-1} X_{i+a,\, j+b}\, K_{a,b}.
\label{eq:q5_conv}
\end{equation}
Sau đó kernel trượt sang một điểm ảnh và lặp lại cho đến khi phủ hết ảnh.
Về mặt toán học, công thức trên là phép tương quan chéo (cross-correlation);
phép tích chập đúng nghĩa cần lật kernel trước khi trượt. Trong học sâu, cả
hai vẫn được gọi chung là tích chập, vì kernel được học từ dữ liệu nên việc
có lật hay không không ảnh hưởng đến khả năng của mô hình. Autograder cũng
dùng phiên bản không lật kernel.

\textbf{Kích thước đầu ra.} Khi không dùng padding và bước trượt (stride)
bằng 1, kernel có thể đặt tại các độ lệch $0, 1, \dots, H - k_H$ theo chiều
dọc và $0, 1, \dots, W - k_W$ theo chiều ngang, nên
\begin{equation}
Y \in \mathbb{R}^{(H - k_H + 1) \times (W - k_W + 1)}.
\label{eq:q5_shape}
\end{equation}
Với ảnh $28 \times 28$ và kernel $3 \times 3$, feature map có kích thước
$26 \times 26$, tức 676 giá trị sau khi duỗi thành vector.

\textbf{Kernel là bộ phát hiện mẫu.} Giá trị $Y_{i,j}$ lớn khi patch tại vị
trí $(i, j)$ có dạng giống với kernel. Ví dụ, kernel
\begin{equation}
K = \begin{bmatrix} 1 & 0 \\ 0 & -1 \end{bmatrix}
\end{equation}
(toán tử Roberts cross) tính hiệu giữa điểm ảnh góc trên trái và góc dưới
phải của mỗi patch $2 \times 2$. Ở vùng có độ sáng đồng đều, hai điểm này
bằng nhau nên đầu ra bằng 0; ở nơi độ sáng thay đổi theo đường chéo, đầu ra
có trị tuyệt đối lớn, và dấu cho biết phía nào sáng hơn. Kernel này vì vậy
phát hiện các cạnh chéo trong ảnh.

\textbf{Tính cục bộ và chia sẻ trọng số.} Mỗi giá trị $Y_{i,j}$ chỉ phụ thuộc
vào một vùng lân cận nhỏ của ảnh (tính cục bộ), và cùng một kernel được dùng
tại mọi vị trí (chia sẻ trọng số). Nhờ đó, một mẫu học được một lần sẽ được
phát hiện ở bất kỳ đâu trong ảnh, đồng thời số tham số rất nhỏ: một kernel
$3 \times 3$ chỉ có 9 tham số, trong khi một lớp tuyến tính
$784 \rightarrow 100$ có $784 \cdot 100 + 100 = 78\,500$ tham số.

\textbf{Kernel được học bằng gradient.} Trong mô hình, kernel là một tham số
và được cập nhật bằng gradient descent giống trọng số của các lớp tuyến
tính. Vì $\partial Y_{i,j} / \partial K_{a,b} = X_{i+a,\, j+b}$, theo quy tắc
chuỗi:
\begin{equation}
\frac{\partial L}{\partial K_{a,b}} = \sum_{i} \sum_{j}
\frac{\partial L}{\partial Y_{i,j}}\, X_{i+a,\, j+b}.
\end{equation}
Nói cách khác, gradient theo kernel chính là phép tương quan chéo giữa ảnh
đầu vào và gradient của feature map. Để PyTorch tự tính được gradient này
(autograd), hàm \texttt{Convolve} phải được viết hoàn toàn bằng các phép
toán tensor.

\textbf{Hàm \texttt{tensordot}.} Lời gọi \texttt{tensordot(a, b, dims=n)}
nhân và lấy tổng theo $n$ trục cuối của \texttt{a} và $n$ trục đầu của
\texttt{b}. Với hai ma trận cùng kích thước và $n = 2$, mọi trục đều bị lấy
tổng nên kết quả là một số vô hướng:
\begin{equation}
\mathrm{tensordot}(A, B, 2) = \sum_{a} \sum_{b} A_{a,b}\, B_{a,b},
\end{equation}
đúng bằng phép ``nhân từng phần tử rồi cộng'' cần cho mỗi vị trí của tích
chập.

Các kiến thức còn lại (MLP với ReLU, softmax và hàm mất mát cross-entropy,
mini-batch SGD) giống Q3 nên không trình bày lại.

% Implement, công thức
\subsubsection{Cài đặt}

\textbf{Hàm \texttt{Convolve}.} Hàm nhận một ảnh $X$ và một kernel $K$, trả
về feature map $Y$ theo các bước:
\begin{enumerate}
  \item Lấy kích thước $H, W$ của ảnh và $k_H, k_W$ của kernel từ thuộc tính
  \texttt{shape}. Kích thước kernel không được cố định là $3 \times 3$, vì
  autograder kiểm tra hàm với kernel $2 \times 2$.
  \item Tính kích thước đầu ra theo công thức~\eqref{eq:q5_shape} và khởi tạo $Y$ bằng hàm
  \texttt{zeros}.
  \item Duyệt mọi vị trí $(i, j)$: cắt patch
  $X_{i:i+k_H,\, j:j+k_W}$ rồi gán $Y_{i,j}$ bằng kết quả
  \texttt{tensordot} của patch với kernel, \texttt{dims=2}.
\end{enumerate}
Kết quả của \texttt{tensordot} được gán trực tiếp dưới dạng tensor, không
chuyển sang số thực Python (ví dụ bằng \texttt{.item()}), vì làm vậy sẽ cắt
đồ thị tính toán và kernel sẽ không nhận được gradient. Với ảnh
$28 \times 28$ và kernel $3 \times 3$, hàm thực hiện 676 lần lặp và
$676 \cdot 9 = 6\,084$ phép nhân cho mỗi ảnh.

\textbf{Kiến trúc.} Mô hình gồm một lớp tích chập với một kernel
$3 \times 3$, theo sau là MLP hai lớp giống Q3:
\begin{align}
Y &= X \star K, & K &\in \mathbb{R}^{3 \times 3},\\
h &= \mathrm{ReLU}(W_1\, \mathrm{vec}(Y) + b_1), & W_1 &\in \mathbb{R}^{128 \times 676},\\
z &= W_2 h + b_2, & W_2 &\in \mathbb{R}^{10 \times 128},
\end{align}
trong đó $\star$ là phép tích chập ở công thức~\eqref{eq:q5_conv} và
$\mathrm{vec}(Y) \in \mathbb{R}^{676}$ là feature map đã duỗi thành vector.
Tổng số tham số là
$9 + (676 \cdot 128 + 128) + (128 \cdot 10 + 10) = 87\,955$, ít hơn mô hình
Q3 ($101\,770$) vì lớp tuyến tính đầu tiên chỉ nhận 676 thay vì 784 đầu vào.

Kernel được khởi tạo bằng ma trận toàn số 1, tương đương một bộ lọc làm mờ
(cộng 9 điểm ảnh lân cận). Mặc dù chú thích trong mã nguồn mô tả mô hình
dùng trọng số huấn luyện sẵn, autograder thực tế gọi hàm huấn luyện, nên
toàn bộ tham số, kể cả kernel, đều được học cùng nhau từ đầu.

\textbf{Huấn luyện.} Mô hình được huấn luyện bằng hàm mất mát cross-entropy
và SGD với learning rate $\eta = 0.01$, batch size $B = 5$, tối đa 50 epoch.
Sau mỗi epoch, độ chính xác trên tập validation được đánh giá, và quá trình
huấn luyện dừng sớm khi đạt 97\%. Ngưỡng này được chọn cao hơn mức 80\% của
autograder và giữ giống Q3 để hai mô hình dễ so sánh.

% Kết quả
\subsubsection{Kết quả}
Mô hình đạt 3/3 điểm autograder. Bảng~\ref{tab:q5_train} cho thấy quá trình
huấn luyện: loss giảm từ 0.87 xuống 0.037 và độ chính xác trên tập validation
tăng từ 84\% lên 98\%, nên huấn luyện dừng sau 7 epoch. Độ chính xác trên tập
test đạt 100\%.
\begin{table}[htbp]
\caption{Quá trình huấn luyện mô hình convolution}
\label{tab:q5_train}
\centering
\begin{tabular}{ccc}
\toprule
Epoch & Loss (batch cuối) & Validation accuracy \\
\midrule
1 & 0.8711 & 0.84 \\
2 & 0.4165 & 0.90 \\
3 & 0.2642 & 0.90 \\
4 & 0.1434 & 0.92 \\
5 & 0.0800 & 0.96 \\
6 & 0.0535 & 0.96 \\
7 & 0.0368 & 0.98 \\
\bottomrule
\end{tabular}
\end{table}

Con số 100\% cần được hiểu đúng. Tập validation và tập test đều được trích
từ 100 ảnh đầu tiên của tập huấn luyện, nên mô hình đã thấy toàn bộ các ảnh
này trong quá trình huấn luyện. Độ chính xác 100\% vì vậy chỉ cho thấy mô
hình khớp tốt với dữ liệu huấn luyện, chứ không đo được khả năng tổng quát
hoá trên ảnh mới. Điều kiện dừng sớm cũng dựa trên dữ liệu đã huấn luyện vì
lý do tương tự. Để đánh giá đúng, cần một tập test tách biệt, ví dụ tập test
MNIST như ở Q3. Đây là hạn chế của cách autograder chia dữ liệu, không phải
của mô hình.

% Hạn chế và mở rộng
\subsubsection{Hạn chế và mở rộng}
\begin{itemize}
  \item \emph{Tốc độ:} hàm \texttt{Convolve} dùng vòng lặp Python với 676 lần
  lặp cho mỗi ảnh, nên chậm hơn nhiều so với các phép toán tensor. Có thể
  vector hoá bằng cách xếp toàn bộ patch thành một tensor kích thước
  $26 \times 26 \times 3 \times 3$, sau đó một lần gọi \texttt{tensordot} với
  kernel (\texttt{dims=2}) cho ra cả feature map $26 \times 26$.
  \item \emph{Chỉ một kernel:} mô hình chỉ phát hiện được một loại mẫu. Các
  CNN thực tế dùng nhiều kernel ở mỗi lớp để tạo nhiều feature map, kết hợp
  padding, stride và pooling, và xếp chồng nhiều lớp tích chập để phát hiện
  các mẫu phức tạp dần, từ cạnh đến nét rồi đến cả chữ số.
\end{itemize}
